% !TEX program = xelatex
% ---------------------------------------------------------------------------
% 摘麦穗与 37% —— 秘书问题与最优停止
% 编译：xelatex "摘麦穗与37%.tex" （跑两遍）
% ---------------------------------------------------------------------------
\documentclass[UTF8,11pt,fontset=macnew]{ctexart}

\usepackage[a4paper,margin=2.4cm]{geometry}
\usepackage{amsmath,amssymb}
\usepackage{booktabs}
\usepackage{array}
\usepackage{enumitem}
\usepackage{xcolor}
\usepackage{tikz}
\usetikzlibrary{arrows.meta,positioning,calc}
\usepackage{pgfplots}
\pgfplotsset{compat=1.18}
\usepackage[most]{tcolorbox}
\usepackage{hyperref}

\definecolor{caramel}{HTML}{A9713C}
\definecolor{rosec}{HTML}{D8678F}
\definecolor{paper}{HTML}{FBF4EA}
\definecolor{inkc}{HTML}{3B2E24}

\hypersetup{colorlinks=true,linkcolor=caramel,urlcolor=rosec,
            pdftitle={摘麦穗与 37\%：秘书问题与最优停止},pdfauthor={米露}}

\linespread{1.13}
\setlength{\parskip}{0.35em}
\setlist{itemsep=0.15em,topsep=0.3em}

\newtcolorbox{keybox}[1][小结]{breakable,enhanced,colback=paper,colframe=caramel,
  boxrule=0.7pt,arc=1.5mm,left=3mm,right=3mm,top=2mm,bottom=2mm,
  title={\bfseries #1},coltitle=white,colbacktitle=caramel,fonttitle=\bfseries}

\newtcolorbox{warnbox}[1][注意]{breakable,enhanced,colback=rosec!6,colframe=rosec,
  boxrule=0.7pt,arc=1.5mm,left=3mm,right=3mm,top=2mm,bottom=2mm,
  title={\bfseries #1},coltitle=white,colbacktitle=rosec,fonttitle=\bfseries}

\title{\bfseries 摘麦穗与 37\%\\[0.35em]
  \large 秘书问题，和一条"看过三成再决定"的法则}
\author{米露的读书笔记}
\date{2026 年 10 月}

\begin{document}
\maketitle
\pagestyle{plain}
\thispagestyle{empty}

\begin{keybox}[这篇文章想做的事]
传说苏格拉底让柏拉图走过一片麦田，不许回头，只许摘一株麦穗，
要摘最大的那一株。柏拉图空着手走了出来。

这个故事的流行版本是一碗励志鸡汤，但它问的其实是一个可以严格数学化的问题：
\textbf{在只能向前走、摘了就不能反悔的前提下，怎样才能最大概率摘到最大的那一株？}
答案是一个很漂亮的数——先扔掉前 $37\%$ 的机会用来校准标准，
之后遇到比前面都好的就摘。成功率约 $1/e\approx36.8\%$。

这篇文章先把它一步步算出来：为什么最优策略长成"观察 + 收割"两段式，
为什么分界点是 $1/e$。然后\textbf{把问题放开两次}：
如果允许摘好几株，出手的节点该落在哪里？
如果把目标从"必须是第一名"换成"名次尽量靠前"，答案又会变成什么样？
（剧透：会变得很不一样——一个给出阈值序列，一个把期望名次压到一个与 $n$ 无关的常数。）
最后才说这个结论在什么条件下会失效。
（本文由大肥鱼写作）
\end{keybox}

\section{麦田里的那条路}

\begin{center}
\begin{tikzpicture}[scale=0.95]
  % 麦田里的行进路线
  \draw[->,thick,draw=caramel] (-0.3,0) -- (10.6,0);
  \foreach \x in {0,0.8,...,10.4} {
    \draw[draw=caramel!70,line width=0.7pt] (\x,0.08) -- (\x,0.55);
    \draw[draw=caramel!70,line width=0.7pt] (\x,0.55) -- (\x+0.14,0.72);
    \draw[draw=caramel!70,line width=0.7pt] (\x,0.55) -- (\x-0.14,0.72);
  }
  % 分界
  \draw[dashed,thick,draw=rosec] (3.7,0) -- (3.7,1.15);
  \node[font=\small,text=rosec,above] at (1.85,1.0) {观察期：一律放弃};
  \node[font=\small,text=caramel,above] at (7.2,1.0) {收割期：遇到"比前面都好"的就摘};
  \node[font=\scriptsize,text=inkc,below] at (0,-0.12) {开始};
  \node[font=\scriptsize,text=inkc,below] at (10.4,-0.12) {走完};
\end{tikzpicture}
\end{center}

故事是这么讲的：苏格拉底把柏拉图带到一片麦田边，让他从一头走到另一头，
不许回头，途中只能摘一株麦穗，而且要摘最大的那一株。
柏拉图一路走一路看，总觉得后面还有更大的，最后两手空空地走了出来。
苏格拉底说：这就是爱情（在另一些版本里是婚姻、是幸福）。

\begin{warnbox}[先说清楚：这个故事是假的]
这段对话\textbf{不见于柏拉图的任何一篇对话录}，也不见于现存的古希腊文献，
是后来（主要在中文的励志读物里）流传开的寓言。
柏拉图笔下的苏格拉底从不这样讲话。

不过——它问的问题是真的，而且有一个非常干净的答案。
\end{warnbox}

要把这个故事变成数学，先得承认一件事：故事里的麦穗是"随机排列"的，
而你\textbf{只能比较你见过的}。你不知道一株麦穗的绝对大小，
只知道它和你之前看过的那些相比是大是小。这一点是整道题的灵魂。

\section{把故事写成模型}

我们把麦田换成一串按随机顺序出现的候选人，于是就有了所谓\textbf{秘书问题}
（secretary problem），也叫择偶问题、相亲问题。

\begin{keybox}[模型与假设]
有 $n$ 个候选者，按\textbf{随机顺序}一个接一个地出现。每见完一个，
你必须立刻决定"要"还是"不要"：

\begin{enumerate}[label=(A\arabic*),leftmargin=2.6em]
  \item 顺序完全随机，$n!$ 种排列等可能；
  \item 只能选一个；一旦选中，过程结束，\textbf{不能反悔、不能回头}；
  \item 你只能判断\textbf{相对}好坏（能比较任意两个见过的候选者），
        不知道任何"绝对"信息（比如分数、身高）；
  \item 目标是让"选中的人恰好是全体最好的那一个"的\textbf{概率最大}——
        选到第二好和选到最差，在目标函数眼里是一样的；
  \item $n$ 已知，而且候选者一定会全部出现（不会中途退出）。
\end{enumerate}
\end{keybox}

第 (A4) 条看着奇怪，但它正是这个模型能被推到底的原因：
目标不是"选个不错的"，而是"全有或全无"。
后面第 \ref{sec:limits} 节会看到，一旦改成"选个够好的就行"，答案就变了。

\section{最优策略长什么样}

一个"策略"说到底就是一个\textbf{停止规则}：在每一刻，根据你看到的东西，
决定是停下还是继续走。下面两个观察极大缩小了候选策略的范围。

\subsection{观察一：非纪录的候选人，永远不该选}

把"到目前为止最好的那个"叫一个\textbf{纪录}（record）。
如果此刻出现的候选人不如前面某个人，那么他\textbf{绝不可能是全体最好的}
（前面那个比他强）。选他就等于放弃。

所以：\textbf{只有纪录才值得考虑}。这一条把"要不要选"的问题
压缩成"要不要选这个纪录"。

\subsection{观察二：值得选的时刻，一定是最后面的一段}

假设你决定在第 $i$ 个纪录出现时接受。那么第 $i+1$ 个纪录出现时，你更应该接受：
剩下的人更少了，机会更稀薄，此时出现的纪录"是全体最好"的把握只会更大。

也就是说，"接受"这个决定的适用范围一定是时间轴上的\textbf{一段尾巴}：
存在一个位置 $r$，在 $r$ 之前无论多好都放弃，从 $r$ 起遇到第一个纪录就选。

\begin{keybox}[最优策略的形状]
存在一个阈值 $r$（$2\le r\le n$），策略是：
\begin{itemize}
  \item 第 $1$ 到第 $r-1$ 个人，\textbf{一律拒绝}（这是观察期，用来建立标准）；
  \item 从第 $r$ 个人起，\textbf{遇到第一个刷新纪录的人就选}。
\end{itemize}
于是整个问题只剩一个未知数：$r$ 取多少最好。
\end{keybox}

（上面的"观察二"是直觉论证。严格版本要用动态规划：设 $V_i$ 为
"第 $i$ 个人是纪录时，从这一刻起继续走下去的成功概率"，
可以证明 $V_i$ 关于 $i$ 单调，因此接受域确实是一段尾巴。
结论不受影响。）

\section{算成功率}

固定阈值 $r$。我们要算的是：这套策略选出全体最优者的概率 $P(r)$。

设最好的那一个恰好出现在第 $i$ 个位置（$i=r,r+1,\dots,n$）。
由随机性，这件事的概率是 $1/n$。

现在假设"最好者在第 $i$ 位"已经成立，要成功还需要什么？
需要你在第 $r$ 位到第 $i-1$ 位之间\textbf{一个纪录都没遇到}——
否则你会提前出手，摘走一个不是最好的。

\begin{keybox}[关键的一步]
在"最好者在第 $i$ 位"的条件下，前 $i-1$ 个人中最好的那一个，
落在 $1,\dots,i-1$ 中任何一个位置的概率都相同。
而"第 $r$ 到第 $i-1$ 位之间没有纪录"恰好等价于
"前 $i-1$ 个人中最强的那个，出现在前 $r-1$ 个人里"。

所以这个条件概率是
\[
\frac{r-1}{i-1}.
\]
\end{keybox}

把两段乘起来再对 $i$ 求和：

\[
\boxed{\;
P(r)\;=\;\sum_{i=r}^{n}\frac{1}{n}\cdot\frac{r-1}{i-1}
\;=\;\frac{r-1}{n}\sum_{j=r-1}^{n-1}\frac{1}{j}\;}
\]

（第二个等号只是把 $j=i-1$ 换了元。）

拿 $n=4$ 验一下：取 $r=2$，则 $P=\frac14\left(\frac11+\frac12+\frac13\right)=0.4583$。
意思是：四个人里\emph{先放弃第一个}，然后遇到比第一个强的就选，
有 $45.8\%$ 的把握选到最好的——比瞎蒙（$25\%$）高出一大截。

\section{连续化：为什么是 \texorpdfstring{$1/e$}{1/e}}

$n$ 大的时候，令 $x=r/n$，把和式看成一个积分：
\[
\sum_{j=r-1}^{n-1}\frac1j\;\approx\;\int_{r-1}^{n}\frac{\mathrm{d}t}{t}\;\approx\;-\ln x,
\]
于是
\[
P(r)\;\approx\;x\,(-\ln x)\;=\;-x\ln x .
\]
对它求导：
\[
\frac{\mathrm{d}}{\mathrm{d}x}\bigl(-x\ln x\bigr)=-\ln x-1,
\]
令它为零得 $\ln x=-1$，即
\[
\boxed{\;x=\frac1e\approx 0.3679\;}
\]
代回去，最大成功率是
\[
-\frac1e\ln\frac1e=\frac1e\approx 36.79\% .
\]

\begin{keybox}[这个 $1/e$ 在说什么]
$P(x)\approx -x\ln x$ 里的两个因子其实就是"观察"与"收割"的
矛盾：

\begin{itemize}
  \item $x$ 是\textbf{观察期}占比。它越大，你的标准越准；
        但 $-\ln x$ 这个因子说明——观察期越长，剩下的收割机会越少
        （$-\ln x$ 随 $x$ 增大而减小）。
  \item 两头都不行：观察期太短（$x\to0$），你还没弄清"什么算好"就出手了；
        观察期太长（$x\to1$），你已经快走完麦田了。
  \item 折中点恰好是 $x=1/e$。
\end{itemize}
\end{keybox}

\begin{center}
\begin{tikzpicture}
\begin{axis}[
  width=12.2cm, height=6.6cm,
  xlabel={$r$（从第几个开始出手）}, ylabel={成功概率 $P(r)$},
  xmin=0, xmax=100, ymin=0, ymax=0.42,
  grid=major, grid style={draw=caramel!18},
  axis line style={draw=caramel!70},
  every axis label/.style={font=\small,text=inkc},
  tick label style={font=\scriptsize},
  legend style={font=\scriptsize,draw=caramel!50,fill=paper,at={(0.98,0.97)},anchor=north east},
]
\addplot[draw=caramel,thick] coordinates {
(1,0.01000) (3,0.08355) (5,0.13376) (7,0.17364) (9,0.20676) (11,0.23484)
(13,0.25890) (15,0.27961) (17,0.29746) (19,0.31281) (21,0.32593) (23,0.33704)
(25,0.34634) (27,0.35397) (29,0.36006) (31,0.36472) (33,0.36804) (35,0.37012)
(37,0.37101) (39,0.37080) (41,0.36953) (43,0.36727) (45,0.36405) (47,0.35992)
(49,0.35492) (51,0.34909) (53,0.34245) (55,0.33505) (57,0.32691) (59,0.31805)
(61,0.30850) (63,0.29829) (65,0.28743) (67,0.27595) (69,0.26386) (71,0.25118)
(73,0.23793) (75,0.22412) (77,0.20978) (79,0.19490) (81,0.17952) (83,0.16363)
(85,0.14726) (87,0.13041) (89,0.11310) (91,0.09533) (93,0.07711) (95,0.05846)
(97,0.03939) (99,0.01990)
};
\addlegendentry{$n=100$ 时的 $P(r)$}
\addplot[draw=rosec,dashed,thick,domain=0:100,samples=2]{0.367879};
\addlegendentry{$1/e\approx0.3679$}
\draw[dashed,draw=rosec] (axis cs:37,0) -- (axis cs:37,0.371);
\end{axis}
\end{tikzpicture}
\end{center}

图里的曲线在 $r\approx37$ 处取到最高点，而且\textbf{非常平}——
这一点很重要，下一节还会回来讲。

\section{精确的最优 \texorpdfstring{$r$}{r}}

连续近似给出 $r\approx n/e$，但真实的最优值是离散的。
用 $P(r+1)-P(r)$ 比较一下相邻两项（推导见本节末）：
\[
P(r+1)-P(r)=\frac1n\left(\sum_{j=r}^{n-1}\frac1j-1\right).
\]
所以 $P(r)$ 上升还是下降，全看这个和式是大于还是小于 $1$。
于是最优阈值是

\[
r^*=\min\left\{r:\ \sum_{j=r}^{n-1}\frac1j\le 1\right\}.
\]

\begin{center}
\begin{tabular}{c c c c}
\toprule
候选人数 $n$ & 最优阈值 $r^*$ & 该拒绝前几个 & 最大成功率 $P(r^*)$ \\
\midrule
3     & 2   & 1    & 0.5000 \\
5     & 3   & 2    & 0.4333 \\
10    & 4   & 3    & 0.3987 \\
20    & 8   & 7    & 0.3842 \\
50    & 19  & 18   & 0.3743 \\
100   & 38  & 37   & 0.3710 \\
1000  & 369 & 368  & 0.3682 \\
10000 & 3680 & 3679 & 0.3679 \\
\midrule
$\to\infty$ & $\approx n/e$ & $\approx 36.8\%$ & $\to 1/e=0.367879$ \\
\bottomrule
\end{tabular}
\end{center}

\begin{warnbox}[一个小细节：$n=100$ 时严格最优是 $r=38$]
按上表，$n=100$ 时最优是"拒绝前 $37$ 个，从第 $38$ 个开始出手"，
对应比例 $38\%$。而 $r=37$ 时的成功率是 $0.371015$，
与最优的 $0.371043$ 只差 $0.00003$——肉眼根本看不出。

所以"$37\%$ 法则"是一个四舍五入的叫法。真正要紧的不是这个数的精确值，
而是它所处的那个平坦区间。
\end{warnbox}

\begin{keybox}[$P(r+1)-P(r)$ 怎么来的]
\begin{align*}
P(r+1)-P(r)
&=\frac1n\left[r\sum_{j=r}^{n-1}\frac1j-(r-1)\sum_{j=r-1}^{n-1}\frac1j\right]\\
&=\frac1n\left[(r-(r-1))\sum_{j=r}^{n-1}\frac1j\;-\;1\right]
=\frac1n\left[\sum_{j=r}^{n-1}\frac1j-1\right].
\end{align*}
（第二步是把第二个和式拆出 $j=r-1$ 那一项，
它正好是 $(r-1)\cdot\frac{1}{r-1}=1$。）
\end{keybox}

\section{这个结论有多结实}

最优解附近太平了，这是 $1/e$ 法则真正的价值所在。
以 $n=100$ 为例：

\begin{center}
\begin{tabular}{c c c}
\toprule
观察期占比 & 相当于拒绝前几个 & 成功率 \\
\midrule
$5\%$  & 5  & $0.134$ \\
$10\%$ & 10 & $0.221$ \\
$20\%$ & 20 & $0.320$ \\
$30\%$ & 30 & $0.363$ \\
$\mathbf{37\%}$ & $\mathbf{37}$ & $\mathbf{0.371}$ \\
$40\%$ & 40 & $0.370$ \\
$50\%$ & 50 & $0.352$ \\
$60\%$ & 60 & $0.313$ \\
$80\%$ & 80 & $0.187$ \\
$100\%$（永远在等）& 100 & $0.010$ \\
\bottomrule
\end{tabular}
\end{center}

\begin{keybox}[读这张表要读出来的两件事]
\textbf{一、$37\%$ 附近很宽容。}在 $30\%$ 到 $50\%$ 之间随便挑一个，
损失都在 $5\%$ 以内。所以你不需要知道 $n$ 的精确值，
只要大致估个量级就行——这对现实很有用。

\textbf{二、两头都很致命。}太早出手（$10\%$）成功率只剩两成；
而"永远觉得后面还有更好的"——走完全程都不出手，
就只能在最后一刻随便摘一株，成功率和瞎蒙一样（$1/n=1\%$）。

所以那个寓言里柏拉图两手空空，从概率上讲其实相当"正常"：
\textbf{这是一场注定大概率失败的游戏}，最优策略也只有 $37\%$。
\end{keybox}

\section{推广一：可以摘多株}
\label{sec:k}

前面一直假设只有一次机会。可现实里我们往往能"多留几个备选"：
相亲可以同时接触几个人，招聘可以发几个 offer，看房子能交几份定金。
如果允许摘 $k$ 株呢？——目标还是老样子：摘到的 $k$ 株里
\textbf{只要有最好的那一株}就算成功。

\subsection{只有纪录值得摘，但阈值不止一个}

"非纪录不可能是最好的"这一条依然成立，所以规则仍然只在纪录上做决定。
新的东西是：\textbf{什么时候开始接受纪录，取决于已经用掉几次机会。}

直觉是这样的：手里已经攒了一株，这次机会的"边际价值"就下降了——
万一第一次摘到的就是全场最好，后面几次根本用不上。
所以用掉一次机会之后，你反而\textbf{更愿意再等等}。

\subsection{递推}

设 $f(i,m)$ 表示：眼下正要查看位置 $i$，已经用掉 $m$ 次机会，
并且已知\textbf{最好的那一株还在 $i,i+1,\dots,n$ 里面}（各位置等可能）时，
最终能摘到那株最好的麦穗的概率。

把"最好的就在眼前"和"最好的还在后面"分开算：

{\small
\[
f(i,m)
=\underbrace{\frac{1}{n-i+1}}_{\text{最好就在位置 }i}\cdot a(i,m)
+\underbrace{\frac{n-i}{n-i+1}}_{\text{最好的在后面}}
\left[\frac1i\cdot\max\bigl(f(i{+}1,m{+}1),\,f(i{+}1,m)\bigr)
+\Bigl(1-\frac1i\Bigr)f(i{+}1,m)\right],
\]}

其中 $a(i,m)$ 在"决定摘"时取 $1$、否则取 $0$；
而"决定摘"当且仅当那个 $\max$ 取在左边（摘比不摘好）。
边界是 $f(n+1,\cdot)=0$；$m=k$ 时不能再摘。

这里用到两件事：位置 $i$ 是纪录的概率是 $1/i$；
以及"$i$ 是不是纪录"与"最好的在后头"相互独立。

算出来的最优规则非常干净：\textbf{对每一个"已用次数" $m$ 都有一个阈值 $t_m$，
位置过了 $t_m$ 之后的纪录就摘，之前的一律放过。} 而且
\[
t_0<t_1<\cdots<t_{k-1},
\]
也就是说：\textbf{用掉的机会越多，出手越晚（越挑剔）。}

\subsection{数值结果}

\begin{center}
\footnotesize
\setlength{\tabcolsep}{7pt}
\begin{tabular}{c c c l}
\toprule
$k$ & $n=100$ 时成功率 & $n\to\infty$ 的极限成功率 & 各阈值的极限位置 $t_0,t_1,\dots,t_{k-1}$ \\
\midrule
1 & 0.3710 & 0.36788 & $0.3679$ \\
2 & 0.5962 & 0.59101 & $0.2231,\ 0.3679$ \\
3 & 0.7385 & 0.73211 & $0.1411,\ 0.2231,\ 0.3679$ \\
4 & 0.8303 & 0.82313 & $0.0910,\ 0.1411,\ 0.2231,\ 0.3679$ \\
5 & 0.8901 & 0.88256 & $0.0594,\ 0.0910,\ 0.1411,\ 0.2231,\ 0.3679$ \\
\bottomrule
\end{tabular}
\end{center}

（阈值一栏是 $t_0,t_1,\dots$：$t_0$ 是第一次出手的时机，最后一个恒为 $1/e$。）

这张表能读出三件事：

\begin{itemize}
  \item \textbf{多一株就多很多。}从 $k=1$ 到 $k=2$，成功率从 $37\%$ 跳到 $60\%$；
        到 $k=5$ 已经接近九成。机会多了，这场游戏就不太像赌博了。
  \item \textbf{最后一个阈值永远等于 $1/e$。}不管一共能摘几株，
        \emph{只剩最后一次}机会时，面对的问题和前面讲的单株问题一模一样，
        所以阈值还是 $0.3679$。
  \item \textbf{前面的阈值明显更早。}$k=2$ 时第一次出手的阈值是 $0.2231$，
        比 $0.3679$ 早了将近一半；$k=5$ 时最靠前那个只有 $0.0594$——
        走过 $6\%$ 之后，见到纪录就可以摘了。
        这就是"有备胎就敢早下手"的数学版本。
\end{itemize}

\begin{keybox}[一个漂亮的巧合]
$0.2231$ 恰好是 $e^{-3/2}$。再往前的那些阈值
（$0.1411$、$0.0910$、$0.0594$）就没有这么简单的表达式了，
只能靠上面的递推翻出来。
\end{keybox}

\section{推广二：换一个目标——让名次尽量靠前}
\label{sec:rank}

到现在为止，目标函数一直是"摘到最好的那株"：
\emph{摘到第二好和摘到最差，在这个目标眼里一模一样}，都算失败。

可现实中大多数人不是这么想的。一个更贴近直觉的目标是：
\textbf{让我摘到的那株尽量靠前}——也就是让它的\textbf{名次}
（$1$ 表示最好）的期望尽量小。

\subsection{目标一换，"只摘纪录"就不够了}

设现在走到位置 $i$，眼前这株在你见过的 $i$ 株里排第 $j$ 名
（$j=1$ 就是纪录）。如果此刻收手，它的名次期望是多少？

由对称性可以算出来：
\[
\text{此刻收手的名次期望}=\frac{j(n+1)}{i+1}.
\]
（理由：你见过的这 $i$ 株在全部 $n$ 株里构成一个均匀随机的 $i$ 元子集，
而这个子集里第 $j$ 小的元素的期望，正是 $\dfrac{j(n+1)}{i+1}$。）

于是决策变成"收手"和"再看一株"比大小：
\[
V(i,j)=\min\Bigl(\underbrace{\frac{j(n+1)}{i+1}}_{\text{现在收手}},\;
\underbrace{C(i)}_{\text{再往后看看}}\Bigr),
\qquad
C(i)=\frac{1}{i+1}\sum_{j'=1}^{i+1}V(i+1,j'),
\]
边界是 $V(n,j)=j$（走到最后一株，必须摘）。

因为 $\frac{j(n+1)}{i+1}$ 随 $j$ 递增，而 $C(i)$ 与 $j$ 无关，
所以最优规则仍然是"看到前几名就摘"：

\begin{keybox}[名次目标下的最优规则]
在位置 $i$，如果当前这株在你见过的 $i$ 株里排进\textbf{前 $J(i)$ 名}，就摘。
\end{keybox}

关键在于：$J(i)$ 现在\textbf{不必等于 $1$}。不像"只要第一名"那个目标，
这里到了后半程，"第三名左右也行了"反而是理性的。
$n=100$ 时算出来的 $J(i)$ 是一条台阶曲线：

\begin{center}
\begin{tikzpicture}
\begin{axis}[
  width=12.2cm, height=5.6cm,
  xlabel={位置 $i$}, ylabel={$J(i)$},
  xmin=0, xmax=100, ymin=0, ymax=54,
  grid=major, grid style={draw=caramel!18},
  axis line style={draw=caramel!70},
  every axis label/.style={font=\small,text=inkc},
  tick label style={font=\scriptsize},
]
\addplot[draw=caramel,thick,const plot] coordinates {
(1,0) (27,0) (28,1) (47,1) (48,2) (59,2) (60,3) (67,3) (68,4) (73,4) (74,5)
(77,5) (78,6) (80,6) (81,7) (82,7) (83,8) (84,8) (85,9) (86,9) (87,10)
(87,10) (88,11) (89,11) (90,12) (91,14) (92,15) (93,17) (94,19) (95,21)
(96,25) (97,30) (98,37) (99,50) (100,50)
};
\draw[dashed,draw=rosec] (axis cs:27,0) -- (axis cs:27,54);
\node[font=\scriptsize,text=rosec,anchor=west] at (axis cs:29,46) {观察期到此结束};
\end{axis}
\end{tikzpicture}
\end{center}

看图要看出两件事：前 $27$ 个位置 $J=0$（一律不摘，纯观察）；
之后 $J$ 一路涨，到最后两个位置时 $J(99)=50$——
意思是"只剩一株了，只要这株比你见过的一半强，就收了吧"。

\subsection{数值结果}

\begin{center}
\begin{tabular}{c c c}
\toprule
$n$ & 最优期望名次 & 首次愿意出手的位置 \\
\midrule
10    & 2.5579 & 第 4 个（$40.0\%$） \\
100   & 3.6032 & 第 28 个（$28.0\%$） \\
1000  & 3.8324 & 第 261 个（$26.1\%$） \\
10000 & 3.8649 & 第 2587 个（$25.9\%$） \\
50000 & 3.8685 & 第 12925 个（$25.9\%$） \\
\midrule
$\to\infty$ & $\to 3.87$ & $\to 25.9\%$ \\
\bottomrule
\end{tabular}
\end{center}

最要紧的是中间那一列：\textbf{最优期望名次趋于一个常数，约 $3.87$。}
换句话说，无论麦田里有多少株麦穗，用这套规则，
你平均摘到的是第 $3.87$ 名上下——\emph{名次不随 $n$ 一起增长}。
这比"成功率三成七"听起来体面多了。

观察期也变了：不是 $37\%$，而是大约 $26\%$。

\subsection{顺带一个提醒：只摘纪录是不够的}

如果把规则限死成"只摘纪录"（也就是前面几节一直在用的那条规则），
$n=100$ 时最好的结果是期望名次 $9.60$（阈值取 $r=10$）；
而允许"摘前几名"的 $J(i)$ 规则能把它压到 $3.60$，好了一倍半还多。

差别出在\textbf{最后一段路}。走到第 $95$ 个位置时，
如果眼前这株是"你见过的第 $21$ 名"，慢条斯理地挑剔似乎很体面；
可你要想一想：不摘它就得继续往前走，而麦田只剩 $5$ 株了。
所以最理性的做法是——\textbf{越接近终点，标准越松}。

\begin{keybox}[两个目标，两种活法]
\begin{center}
\small
\begin{tabular}{@{}l l l@{}}
\toprule
目标 & 最优规则 & 结果 \\
\midrule
"一定要摘到最好的" & 跳过前 $37\%$，之后见到纪录就摘 & 成功率 $1/e\approx36.8\%$ \\
"名次尽量靠前" & 跳过前 $26\%$，之后排进前 $J(i)$ 名就摘 & 期望名次 $\approx3.87$ \\
\bottomrule
\end{tabular}
\end{center}

同一片麦田，目标只改了一句话，最优点就从 $37\%$ 挪到了 $26\%$，
规则也从"只认第一名"松成"越晚越将就"。
\textbf{最优停止对目标函数极其敏感}——这是这两节最值得记的一句话。
\end{keybox}

\section{两种推广合起来}

如果既能摘 $k$ 株，目标又是"摘到的最好那株名次尽量靠前"呢？
这时候要看的是\emph{摘到的 $k$ 株里最好的那一株}的名次。
用同样的动态规划可以算，只是状态里要多记一个
"已经摘到的最优者前面还有几株比它好"。

\begin{center}
\begin{tabular}{c c c}
\toprule
机会数 $k$ & $n=100$ 时的期望名次 & $n=200$ 时的期望名次 \\
\midrule
1 & 3.6032 & 3.7192 \\
2 & 1.9573 & 1.9812 \\
3 & 1.4953 & 1.5085 \\
4 & 1.2818 & 1.2914 \\
5 & 1.1671 & 1.1748 \\
\bottomrule
\end{tabular}
\end{center}

多几次机会，期望名次掉得很快，但\textbf{掉到 $1$ 附近就停住了}——
名次不可能比 $1$ 更小。$k=5$ 时期望名次 $1.17$，
大致意思是：多数时候你都能摘到前三名。

\section{什么时候这个结论会失效}
\label{sec:limits}

$1/e$ 法则依赖前面那五条假设，缺一条答案就会变。这些假设在现实生活中
几乎从来不成立，所以把它当人生教条用是危险的：

\begin{center}
\begin{tabular}{p{5.6cm} p{8.2cm}}
\toprule
假设 & 现实中往往是什么样 \\
\midrule
(A1) 顺序随机 & 候选人常常扎堆出现，好坏有相关性；"今年行情好"这种信息会让模型失效 \\
(A2) 不能反悔 & 现实中常常可以回头联系，或者先"放在池子里"；一旦允许回头，
                最优策略会完全不同（而且成功率会高很多） \\
(A3) 只有相对信息 & 我们通常知道一些绝对值（学历、薪资、相处的感觉）；
                有绝对信息时，问题变成"全信息最优停止"，策略和答案都不一样 \\
(A4) 目标是"选到最好的" & 大多数人要的是"足够好"，不是"第一名"。
                第 \ref{sec:rank} 节算过：只把目标换成"名次尽量靠前"，
                观察期就从 $37\%$ 缩到 $26\%$，规则也从"只摘纪录"松成"越晚越将就"。
                至于"只摘一株"这条，第 \ref{sec:k} 节也放开了 \\
(A5) $n$ 已知 & 现实中你通常不知道这辈子会遇到多少人；
                所幸前面说过，$30\%\sim50\%$ 这片区域很平，
                $n$ 估错一点问题不大 \\
\bottomrule
\end{tabular}
\end{center}

还有一种更彻底的改法：给每位候选者一个\textbf{绝对分数}，而不只是相对名次。
这时目标可以写成"让被选中者的期望分数最大"，
而答案会变得完全不同——最优策略是一串\textbf{随时间下降的分数线}，
而不是"要不要刷新纪录"。这套东西常被叫做\textbf{全信息问题}
（full-information problem）；它和第 \ref{sec:rank} 节的"名次期望"问题
只差一点措辞，答案却差得很远。

\begin{keybox}[一句话总结这些变体]
$1/e$ 法则是"在最极端的无知下仍然能做到的最好"。
它之所以值得记住，不是因为 $37\%$ 这个数，
而是因为它揭示了一个结构：\textbf{任何"边看边决定"的问题里，
都存在一个"先用一部分机会买信息、再用剩下的机会收割"的权衡}。
这个结构在最优停止理论里到处出现，秘书问题只是它最干净的一个特例。
\end{keybox}

\section{小结}

\begin{enumerate}
  \item 把麦田问题写成模型：$n$ 个候选者随机顺序出现，
        只能即时决定、不能回头，目标是选到最好的那个。
  \item 只有"纪录"值得考虑；而且值得接受的时刻一定是时间轴末尾的一段。
        所以最优策略只需一个参数 $r$：\textbf{前 $r-1$ 个一律放弃，
        之后遇到第一个纪录就选}。
  \item 成功率
        $P(r)=\dfrac{r-1}{n}\displaystyle\sum_{j=r-1}^{n-1}\dfrac1j$，
        连续化后为 $P\approx -x\ln x$（$x=r/n$），
        在 $x=1/e$ 处取最大值 $1/e\approx36.8\%$。
  \item 精确最优阈值是满足 $\sum_{j=r}^{n-1}\frac1j\le1$ 的最小 $r$，
        约为 $n/e$。
  \item 结论在 $30\%\sim50\%$ 这一大片区间里都很稳，但"太早"和"一直等"
        都会崩塌。
  \item 允许摘 $k$ 株（第 \ref{sec:k} 节）：最优策略是一串阈值
        $t_0<t_1<\cdots<t_{k-1}$——用掉几次机会之后反而更挑剔；
        最后一个阈值恒为 $1/e$。$k=2$ 就能把成功率从 $37\%$ 提到 $59\%$。
  \item 换成"名次尽量靠前"（第 \ref{sec:rank} 节）：规则松成
        "排进前 $J(i)$ 名就摘"，观察期缩到 $26\%$，
        而\textbf{期望名次收敛到一个与 $n$ 无关的常数 $\approx3.87$}。
\end{enumerate}

故事里的苏格拉底并没有给出策略，他只留下了一句"这就是爱情"。
数学没能让这件事变得容易，但至少把它变清楚了：
\textbf{三十六点八，且不能回头。}
——当然，这句话只在"非摘到第一名不可"的那个版本里成立。
把目标换一换，数字就跟着变；
这也正是最优停止这件事最容易被鸡汤化的地方。

\begin{keybox}[参考]
\begin{itemize}
  \item 秘书问题的标准综述与完整推导，见
        T.~S.~Ferguson, \textit{Who Solved the Secretary Problem?},
        Statistical Science \textbf{4} (1989), 282--289；
        以及 P.~R.~Freeman, \textit{The Secretary Problem and its Extensions:
        A Review}, International Statistical Review \textbf{51} (1983), 189--206。
  \item 名次期望（本文第 \ref{sec:rank} 节）这条路线最早的系统研究之一，见
        Y.~S.~Chow, S.~Moriguti, H.~Robbins, S.~M.~Samuels,
        \textit{Optimal selection based on relative rank
        (the ``secretary problem'')}, Israel Journal of Mathematics
        \textbf{2} (1964), 81--90；
        极限期望名次 $3.8695\ldots$ 就出自这条路线。
  \item 可以摘 $k$ 株（本文第 \ref{sec:k} 节）的完整处理见
        J.~P.~Gilbert, F.~Mosteller, \textit{Recognizing the maximum of a
        sequence}, Journal of the American Statistical Association
        \textbf{61} (1966), 35--73。
  \item 面向大众的讲法，见 B.~Christian, T.~Griffiths,
        \textit{Algorithms to Live By}（中译本《算法之美》），
        其中一章专讲最优停止。
\end{itemize}
\end{keybox}

\end{document}
